\section{Limitations}
\label{sec:limitations}

The experiments concern cross-condition transfer with bearing-disjoint targets on two test rigs; no cross-machine claim is
made, and the two rigs play different roles --- Paderborn carries the full analysis, HUST only the replication of the
collapse. The HUST splits change bearing size together with bearing identity, so they confirm the effect without isolating
its cause, and HUST's unpublished pitch diameter keeps it out of every kinematic analysis.

Bearing-disjoint splits on Paderborn move more than bearing identity. Among the real-damage outer-race bearings, KA04, KA16
and KA22 carry fatigue pitting while KA15 and KA30 carry plastic deformation from indentation, so a split can place one
damage mechanism mostly on one side. The design controls the operating condition exactly, not the mechanism or the
difficulty of the individual specimen; ten random splits average over that variation but do not remove it, and a
same-condition control that varies only the bearing set would isolate it further.

Seeds and splits are traded against each other: three seeds on the two mechanical folds, one seed on each of the ten random
splits and six HUST splits. The paired design makes this affordable, since the quantity of interest is a within-job
difference, but a full seeds-within-splits grid was outside the compute budget and would tighten the split-level intervals.

Statistically, the two mechanical folds provide only two clusters, so their task-level $p$-values are descriptive and the
inference rests on the ten random splits (Section~
ef{sec:protocol}). The identity probes are measured on held-out
recordings of bearings the source model was trained on, so they establish what the representation encodes, not how it
behaves on unseen specimens.

Two source-free methods are evaluated. Neighbourhood-based methods such as NRC and AaD, and 2026-generation bearing-specific
methods, remain to be tested under this protocol; nothing here predicts that they behave differently, but nothing here
demonstrates it either.

The gating results inherit the coverage limit of envelope analysis on this dataset: the envelope gate engages 7 of the 11
real-damage bearings and fires on about a quarter of their recordings, ball faults are absent from the Paderborn real-damage label space, and no gate we tested detects rolling-element
faults on CWRU or UORED. Slip was measurable on \SI{16}{\percent} of 6203 fault recordings, namely those with at least two
lit harmonics, and UORED is excluded from the slip analysis for lack of an independent speed reference.

Finally, the raw-spectrum band rule is our implementation of a published description, not the authors' code, which was not
available; we therefore report it as a mechanism and never compare it against the accuracy published for the method it comes
from. A remedy that removes bearing identity from the source representation --- for example a class-conditional adversary,
which the probe of Section~\ref{sec:res-probe} motivates directly --- was pre-registered but cut for budget before any
outcome was observed, and is left as the obvious next experiment.
